\documentclass{amsart}
\usepackage[T1]{fontenc}
\usepackage{amssymb}
\usepackage[hidelinks]{hyperref}
\newtheorem{theorem}{Theorem}
\newtheorem{lemma}{Lemma}[section]
\newtheorem{fact}{Fact}
\renewcommand{\thefact}{F\arabic{fact}}
\theoremstyle{remark}
\newtheorem*{remark}{Remark}
\newcommand{\rad}{\operatorname{rad}}
\newcommand{\Q}{\mathbb Q}
\newcommand{\Z}{\mathbb Z}
\newcommand{\eps}{\varepsilon}
\numberwithin{equation}{section}
\title{Improving the Stewart--Yu bound for all abc triples}
\hypersetup{pdftitle={Improving the Stewart--Yu bound for all abc triples}}
\author{Ruiyi Zhang}
\email{ryzfsu@gmail.com}
\address{Independent Researcher, Columbia, Maryland, USA}
\hypersetup{pdfauthor={Ruiyi Zhang}}
\subjclass[2020]{Primary 11D75; Secondary 11J86, 11D61}
\keywords{abc conjecture, radical, effective bounds, $p$-adic logarithmic forms}
\date{November 13, 2026}
\begin{document}
\begin{abstract}
We prove that $\log c<A R^{1/3}\exp(-\log R/(3000\log\log R))$
for coprime positive integers $a+b=c$, where $R=\operatorname{rad}(abc)$
and $A$ is absolute and effectively computable.
The proof combines estimates at maximal prime powers with a growing number
of auxiliary primes in a rational logarithmic form.
\end{abstract}
\maketitle

\section{Introduction}
For positive coprime integers $a+b=c$, put $R=\rad(abc)$.
The $abc$ conjecture predicts $c<K_\eps R^{1+\eps}$ for every $\eps>0$.
It remains open.
Stewart--Yu \cite[p.~226]{SY1991} proved
$\log c<R^{2/3+C/\log\log R}$ for $c>2$.
Their later bound \cite[Theorem~1, p.~170]{SY2001} is
$\log c\ll R^{1/3}(\log R)^3$, effectively.
We prove the following.

\begin{theorem}\label{thm:main}
There is an absolute effectively computable $A>0$ such that
\begin{equation}\label{eq:main}
\log c<A R^{1/3}\exp\!\left(-\frac{\log R}{3000\log\log R}\right)
\end{equation}
for all positive integers $a+b=c$ with $\gcd(a,b)=1$.
\end{theorem}
Equivalently, $\log c<A R^{1/3-1/(3000\log\log R)}$, where the exponent tends to $1/3$ from below as $R\to\infty$.
Except for $(1,1,2)$, one has $R\ge6$ and $\log\log R>0$.
In particular, $\log c\ll R^{1/3}$ effectively.
For $R\ge6$, Pasten \cite[Theorem~1.4]{Pasten2023arXiv} bounds $\log c$ by
$P_{\min}\exp(\kappa_{\rm P}\sqrt{\log R\log\log R})$.
Here $P_{\min}=\min\{P(a),P(b),P(c)\}$, $P(n)$ is the largest prime factor and $P(1)=1$.
For $a\le c^{1-\vartheta}$ and fixed $\vartheta>0$, he bounds it by
$\vartheta^{-1}\exp(\kappa_{\rm P}\sqrt{\log R\log\log R})$,
with an absolute $\kappa_{\rm P}>0$.
These improve the bound in restricted regions.
Neither displayed estimate gives an improvement over the Stewart--Yu bound $\log c\ll R^{1/3}(\log R)^3$ for all triples in terms of $R$ alone.

Put
\begin{equation}\label{eq:notation}
H=\log c,\qquad W=\log R,\qquad L=1+W,\qquad X=R^{1/3}.
\end{equation}
Let $M_i,s_i,\Delta$ be as in Section~\ref{sec:complement}.
Let $C_T$ be the effective constant in \eqref{eq:outerconstants} and put
$C_P=100000C_V$, with $C_V$ defined in \eqref{eq:CV}.
On $\mathcal N$ in \eqref{eq:N}, we use
\begin{equation}\label{eq:skeleton}
H<C_T XL^2\Delta^{-2/9},\qquad
H<C_P XL^{10}\exp\!\left(-\frac{W}{156\log W}
                         +\frac{11}3\log\Delta\right).
\end{equation}
The second holds for $W\ge W_0$ when $\log\Delta\le(3/1820)W/\log W$.
At this boundary both rates equal $1/2730>1/3000$.
Section~\ref{sec:proof} gives an effective common threshold.
Stewart \cite[Section~4, (4.13)--(4.14)]{Stewart2013} introduced the
growing auxiliary-prime construction.
Section~\ref{sec:cap} bounds the cofactors $u_i$ used in applying this
construction to additive triples.
For earlier $p$-adic estimates, see \cite{Yu2007}.
The use of AI tools is described at the end of the paper.

\section{Known estimates}\label{sec:inputs}
All logarithms are natural.
Write $\omega(n)$ for the number of distinct prime divisors of $n$,
with $\omega(1)=0$, and $\pi(x)$ for the number of primes at most $x$.
For reduced $x/y$, put $h(x/y)=\log\max(|x|,|y|)$.
Write $T(z)=\max(1,h(z))$ for rational $z$.
The valuation satisfies $v_p(p)=1$ and $|x|_p=p^{-v_p(x)}$.

\begin{fact}[{\cite[Theorem~1, p.~170, (5)]{SY2001}}]\label{fact:global}
There is an absolute effectively computable $S\ge1$ such that
\begin{equation}\label{eq:global}
H\le S XW^3.
\end{equation}
\end{fact}

\begin{fact}[{\cite[Lemma~2, p.~173, and (10), p.~172]{SY2001}}]\label{fact:local}
For rational $p$-units $\alpha_1,\ldots,\alpha_n$, integer exponents
$|b_i|\le B_0$, $B_0\ge3$, and $\Theta=\prod_i\alpha_i^{b_i}-1\ne0$,
\begin{equation}\label{eq:local}
v_p(\Theta)<\frac p{\log p}(c_9n)^{2n}
 \prod_i\frac{\max(h(\alpha_i),\log p)}{\log p}\log B_0,
\end{equation}
where $c_9>0$ is absolute and effectively computable.
\end{fact}
In $K=\Q(\alpha_0)$ of \cite[Section~2, (10)]{SY2001}, the ramification
index at $p$ is $1$, including $p=2$; hence $\operatorname{ord}_{\mathfrak p}$
on $\Q$ agrees with $v_p$.
Independence is not required, and a base $-1$ is allowed.

\begin{fact}[{\cite[Proposition~5, p.~517]{Gyory2019}, corrected in
\cite{Gyory2020}}]\label{fact:group}
Put $K_0=300(16e)^8$.
Let $m\ge1$ and let $\Gamma$ be generated by distinct primes
$q_1,\ldots,q_m$ and $-1$.
The primes form a basis for $\Gamma/\Gamma_{\rm tors}$.
If $\xi\in\Gamma$ and $\xi\ne1,-1$, then
\begin{align}
-\log|1-\xi|&<K_0^m\log\max(e,h(\xi))\prod_j\log q_j,\label{eq:real}\\
-\log|1-\xi|_p&<K_0^m\frac p{\log p}
 \log\max(e,ph(\xi))\prod_j\log q_j.\label{eq:finite}
\end{align}
\end{fact}
\begin{lemma}[Specialisation of the group estimate]\label{lem:group-specialisation}
Proposition~5 of \cite{Gyory2019}, with the sign corrected in \cite{Gyory2020},
gives \eqref{eq:real}--\eqref{eq:finite}
under the hypotheses of Fact~\ref{fact:group}.
\end{lemma}
\begin{proof}
Apply Proposition~5 of \cite{Gyory2019}, with the sign corrected in \cite{Gyory2020},
with $K=\Q$ and the parameter $\alpha=1$ in \cite{Gyory2019}.
The group has positive rank, and $h(\xi)\ge\log2>0$.
Its constant is $48(16e)^8$ for $m=1$.
For $m\ge2$, the specialized coefficient
$2(m+1)\log\max(e,m)(16e)^{3m+5}$ is at most
$4^m(16e)^{8m}\le\bigl(8(16e)^8\bigr)^m$, since
$m+1\le2^m$, $\log\max(e,m)\le2^{m-1}$ and $3m+5\le8m$.
The factor for the real place is at most $F=4/\log2>1$.
For $m\ge2$, $F\bigl(8(16e)^8\bigr)^m
\le\bigl((32/\log2)(16e)^8\bigr)^m<K_0^m$.
For $m=1$, $48F<300$; both comparisons follow from $\log2>2/3$.
\end{proof}

\begin{fact}[{\cite[Lemma~3.1, pp.~301--302]{Stewart2013};
\cite[Main theorem, (1.18), p.~321]{Yu2013}}]\label{fact:growing}
Let $p\ge5$, $n\ge7$, and $\ell=\log p$.
For independent rational $p$-units $\alpha_i$ and integers $b_i$, not all
zero, put $B_e=\max(2,|b_1|,\ldots,|b_n|)$.
If $\prod_i\alpha_i^{b_i}\ne1$, then
\begin{equation}\label{eq:growing}
\begin{split}
v_p\left(\prod_i\alpha_i^{b_i}-1\right)
<&\ 376\sqrt{n+1}\,[4+\log(n+1)]
 \left(7e\frac{p-1}{p-2}\right)^n\prod_i h(\alpha_i)\\
&\cdot\max\left(p(n/\ell)^n,e^n\ell\right)
 \max\bigl(\log B_e,5.4n(n+1)\bigr).
\end{split}
\end{equation}
\end{fact}
\begin{lemma}[Specialisation of the growing-rank estimate]\label{lem:growing-specialisation}
Under the hypotheses of Fact~\ref{fact:growing}, estimate \eqref{eq:growing}
follows from its cited sources.
\end{lemma}
\begin{proof}
Stewart's derivation uses the constants $1794$ and $5.25$.
Yu's published rational case has $1790$ and $5.84$; with these constants,
the bound $G_1\le5.4n(n+1)$ used here holds for $n\ge7$.
We derive this range directly from Yu~\cite{Yu2013}.
Use \cite[(1.9), (1.11), (1.13), (1.28), (1.30)--(1.31)]{Yu2013}
with $d=f=e_{\mathfrak p}=1$ and Yu's parameters
$q_{\rm Y}=2$, $q_{\rm Y}^{u_{\rm Y}}=2$.
Relabel a minimum-valuation exponent last, as (1.17) requires.
We retain $5.84$ throughout and use the upper bound $1790\le1794$.
Put $Q_n=4+\log(n+1)$ and $M_e=\max(\log B_e,5.4n(n+1))$.
Then
\[
\frac{G_1}{n+1}<4n+6+\log(4n+6)\le5.4n\quad(n\ge7).
\]
The last inequality holds at $7$, and its difference increases.
Since $h(\alpha_i)\ge\log2$ and the auxiliary quantity in
\cite[(1.15)]{Yu2013} equals $\log2\log3/\log6$,
the first term of $h^{(1)}$ is at most
\[
\log B_e+\log(2\log3/\log6)<\log B_e+\tfrac14<1.001M_e.
\]
Yu's additional term $\log B^\circ$, where
$B^\circ=\min\{|b_i|:b_i\ne0\}$, satisfies
$\log B^\circ\le\log B_e\le M_e$.
Hence $h^{(1)}\le\max(1.001M_e,(n+1)\ell)$.
The saturated residue-group index $\delta(\boldsymbol\alpha)$ is at least $1$.
Yu's main theorem uses the heights of the original independent tuple
in both Kummer cases; we only enlarge $p/\delta(\boldsymbol\alpha)$ to $p$.
Yu's first bound is consequently at most
\[
\begin{split}
&1794\left(7\frac{p-1}{p-2}\right)^n
 \frac{(n+1)^{n+1}}{2n!\ell^2}\prod_i h(\alpha_i)\\
&\quad\cdot\max\left(p(n/\ell)^n,e^n\ell\right)
 \max(Q_n,\ell)\max(1.001M_e,(n+1)\ell).
\end{split}
\]
Use Stirling's inequality
$(n+1)^{n+1}/n!\le e^{n+1}\sqrt{n+1}/\sqrt{2\pi}$.
Also $\max(Q_n,\ell)/\ell\le Q_n/\log5$ and
$\max(1.001M_e,(n+1)\ell)/\ell\le1.001M_e/\log5$.
Finally,
\[
\frac{1.001\cdot1794e}{2\sqrt{2\pi}(\log5)^2}<376.
\]
For this comparison one may use
$e<2.718282$, $\log5>1.609437$, and $\sqrt{2\pi}>2.50662$.
This gives \eqref{eq:growing} from the published constants for $n\ge7$.
\end{proof}

Put $C_H=6+\log S$.
By \eqref{eq:global}, if $h(\xi)\le H$ and $1\le y\le R$, then
\begin{equation}\label{eq:logheight}
\log\max(e,yh(\xi))\le C_HL,
\end{equation}
since $\log(RH)\le\log S+4W/3+3\log W\le C_HL$ and $C_HL\ge6$.
For fixed $\gamma>0$, define the effective constants
\begin{equation}\label{eq:calpha}
\begin{split}
U_\gamma&=\left\lceil\max(1,4/\gamma^2,2\log K_0/\gamma)\right\rceil,\\
C_\gamma&=\prod_{\substack{q<e^{U_\gamma}\\q\ \mathrm{prime}}}
 \max\left(1,\frac{K_0\max(1,\log q)}{q^\gamma}\right).
\end{split}
\end{equation}
For every squarefree $J$,
\begin{equation}\label{eq:absorbprimes}
G(J):=K_0^{\omega(J)}\prod_{q\mid J}\log q\le C_\gamma J^\gamma,
\qquad G(1)=1.
\end{equation}
For $\zeta=\log q\ge U_\gamma$,
$\log K_0\le\gamma\zeta/2$ and $\log\zeta\le\sqrt\zeta\le\gamma\zeta/2$.
Multiply over the primes dividing $J$.

\section{Maximal prime powers}\label{sec:complement}
Arrange $a\le b$.
For each nonunit side $i\in\{a,b,c\}$, put
\[
M_i=\max_{q\mid i}v_q(i),\qquad
s_i=\min\{q:v_q(i)=M_i\},\qquad
u_i=i/s_i^{M_i},\qquad E_i=\rad(u_i).
\]
By definition, $s_i\nmid u_i$.
On $a^2>c$, all three sides exceed $\sqrt c$.
Put
\[
x_i=\log s_i,\quad H_i=M_ix_i,\quad
\Delta=E_aE_bE_c,\quad J_i=\Delta/E_i,\quad \Omega=\log\Delta.
\]
Consider the range
\begin{equation}\label{eq:N}
\mathcal N=\{a^2>c,\ s_i^{39}>R^{10}\text{ for every }i\}.
\end{equation}
On this range,
\begin{gather}
R=s_as_bs_c\Delta,\qquad J_aJ_bJ_c=\Delta^2,\label{eq:identities}\\
\frac{10}{39}W<x_i<\frac{19}{39}W,\quad
\Omega<\frac9{39}W<x_i,\quad
\log\rad(i)<\frac{19}{39}W,\quad H_i>\frac5{19}H.\label{eq:comparisons}
\end{gather}
We have $x_i\le W-x_j-x_k<19W/39$ and
$\log\rad(i)\le W-x_j-x_k$; summing the three lower bounds gives
$\Omega=W-x_a-x_b-x_c<9W/39$.
Maximality gives
$H_i\ge x_i\log i/\log\rad(i)>(10/19)\log i>5H/19$.
Also $R\ge30$, $W>3$ and $W>3L/4$.
The actual ratios
\begin{equation}\label{eq:ratios}
\xi_a=b/c,\qquad\xi_b=a/c,\qquad\xi_c=-a/b
\end{equation}
have height at most $H$ and satisfy
$v_q(1-\xi_i)=v_q(i)$ for $q\mid i$.

Set
\begin{equation}\label{eq:outerconstants}
C_U=2C_HC_{1/6},\qquad
C_S=\frac{2C_HK_0^2C_{1/26}}{(\log2)^2},\qquad
C_T=\frac{741}{50}C_HK_0^2C_{1/6}.
\end{equation}
If $a^2\le c$, apply (F3) at infinity to $\xi=b/c$.
Here $-\log(1-\xi)\ge H/2$.
Equations \eqref{eq:logheight}--\eqref{eq:absorbprimes} give
\begin{equation}\label{eq:U}
H<C_U XL e^{-W/6}.
\end{equation}
This includes unit sides and $a^2=c$.
If $a^2>c$ and some $s_i^{39}\le R^{10}$, choose
$p=s_i=\min\{s_a,s_b,s_c\}$; then $p^{39}\le R^{10}$.
Since $\log i\le M_i\log\rad(i)\le M_iW$ and $\log i>H/2$,
maximality gives $H<2WH_i/\log p$.
Apply (F3) to the full off-side prime basis in \eqref{eq:ratios}:
\[
H_i<C_HK_0^2C_{1/26}\frac p{\log p}Lx_jx_kJ_i^{1/26}
 \le C_HK_0^2C_{1/26}\frac p{\log p}L^3R^{1/26}.
\]
Since $p\le R^{10/39}=X R^{-1/13}$,
\begin{equation}\label{eq:S}
H<C_S XL^4e^{-W/26}.
\end{equation}

On all of $\mathcal N$, (F3) together with \eqref{eq:absorbprimes},
with $\gamma=1/6$, gives
\[
\begin{gathered}
\frac{x_jx_k}{x_i}\le\frac{39}{10}W,\qquad
C_*=\frac{39}{10}C_HK_0^2C_{1/6},\\
H_i<C_HK_0^2C_{1/6}L\frac{s_i}{x_i}x_jx_kJ_i^{1/6}
 <C_*s_iL^2J_i^{1/6}.
\end{gathered}
\]
Multiply and use \eqref{eq:identities}--\eqref{eq:comparisons}:
\begin{equation}\label{eq:thick}
\prod_iH_i<C_*^3RL^6\Delta^{-2/3},\qquad
H<C_T XL^2\Delta^{-2/9}.
\end{equation}
There is no restriction on $\Delta$ in this estimate.

\section{Cofactors of the maximal prime powers}\label{sec:cap}
Assume the triple lies in $\mathcal N$.
Put $\kappa=\max(1,c_9)$ and $F(n)=(\kappa n)^{2n}$.
Choose an integer $m_0\ge7$ with $\log(m_0+1)\ge8\log\kappa+10$.
Define
\begin{equation}\label{eq:CV}
C_F=\max\left(1,\max_{0\le m<m_0}
 \frac{F(m+3)}{((m+1)!)^3}\right),\qquad
C_V=\frac{(C_H+2)C_F}{4(\log2)^3}.
\end{equation}
Then $C_V\ge1$, since $C_F\ge F(3)\ge729$.
For squarefree $J$ with $m=\omega(J)$, one has $J\ge(m+1)!$.
Integration gives $\log((m+1)!)\ge(m+1)\log(m+1)-m$.
It follows that
\[
\begin{split}
\log\frac{F(m+3)}{((m+1)!)^3}
&\le(3-m)\log(m+1)+2(m+3)\log\kappa+3m+4+\frac8{m+1}\\
&\le m\left[-\tfrac12\log(m+1)+4\log\kappa+5\right]\quad(m\ge7).
\end{split}
\]
The last expression is nonpositive for $m\ge m_0$.
The finite maximum handles $m<m_0$, including $0$.
Consequently,
\begin{equation}\label{eq:factorial}
F(3+\omega(J))\le C_F J^3.
\end{equation}

Let $l\mid E_i$.
Represent $\xi_i$ using $s_j,s_k$, every prime dividing $J_i$, and the sign.
These bases are $l$-units.
Their coefficients are at most $B_0=\max(3,H/\log2)$, with
$\log B_0\le(C_H+2)L$.
Since $\log l\le\Omega<x_j,x_k$, the two factors for $s_j,s_k$ in
\eqref{eq:local} have product at most $L^2/[4(\log2)^2]$.
Each remaining prime contributes
\[
\max(1,\log q/\log l)\le\max(1,\log q/\log2)<q.
\]
The sign contributes $1$.
Apply (F2), \eqref{eq:factorial} and $l/\log l\le E_i/\log2$:
\begin{equation}\label{eq:residual}
1\le v_l(i)<C_VL^3E_iJ_i^4\le C_VL^3\Delta^4.
\end{equation}
Every residual exponent is included, also at $l=2$.
For distinct sides $j,k$, this yields
\begin{equation}\label{eq:cofactor}
T(u_j/u_k)\le C_VL^4\Delta^4.
\end{equation}
If $u_j>1$,
$\log u_j\le\max_{l\mid u_j}v_l(j)\,\Omega$ and $\Omega<W<L$;
if $u_j=1$, $\log u_j=0$.
Likewise for $u_k$. Coprimality gives
$h(u_j/u_k)=\max(\log u_j,\log u_k)$; use $C_V\ge1$ for $T(u_j/u_k)$.

\section{Auxiliary primes in the local estimate}\label{sec:aux}
\begin{lemma}\label{lem:aux}
Let $p,q,r$ be distinct primes, $\ell=\log p$, and $\Lambda=\log\ell\ge10^6$.
Suppose $\ell\le\log q,\log r\le2\ell$.
Let $w_q,w_r$ be coprime positive integers with
$\gcd(w_qw_r,pqr)=1$ and $\omega(w_qw_r)\le\ell$.
For $1\le N_q,N_r\le p^2$, $\sigma\in\{1,-1\}$, and a nonzero difference,
\begin{equation}\label{eq:auxbound}
v_p(w_qq^{N_q}-\sigma w_rr^{N_r})
 <25200p\ell^5T(w_q/w_r)
                     \exp\!\left(-\frac{\ell}{52\log\ell}\right).
\end{equation}
\end{lemma}
\begin{proof}
Write $\beta$ for the valuation and $z=(w_q/w_r)q^{N_q}r^{-N_r}>0$.
If $\beta=0$, the result follows.
For $\sigma=1$, take $\nu=1$.
For $\sigma=-1$, take $\nu=2$; oddness gives
$v_p(z^2-1)=v_p(z+1)=\beta$.
Hence $\beta=v_p(z^\nu-1)$.
The denominator is a $p$-unit, and $v_q(z)=N_q>0$ ensures nonvanishing.

Put $n=\lfloor\ell/(52\Lambda)\rfloor$.
Then $n\ge7$ and $n\ge\ell/(104\Lambda)$.
Choose auxiliary primes outside $pqrw_qw_r$.
For $\theta=w_q/w_r\ne1$, use $n-3$ of them and set
\begin{equation}\label{eq:tuple}
\begin{aligned}
\alpha_1&=q/(\varpi_4\cdots \varpi_n),& b_1&=\nu N_q,\\
\alpha_2&=1/r,& b_2&=\nu N_r,\\
\alpha_3&=\theta,& b_3&=\nu,\\
\alpha_i&=\varpi_i\ (i\ge4),& b_i&=\nu N_q.
\end{aligned}
\end{equation}
If $w_q=w_r=1$, omit $\theta$ and use $n-2$ auxiliaries.
The product is $z^\nu$.
Take an integer relation $\prod_i\alpha_i^{z_i}=1$.
The valuations at $q$ and $r$ give $z_1=z_2=0$.
For $\theta\ne1$, a prime in its support then gives $z_3=0$,
and each $\varpi_i$ gives $z_i=0$ for $i\ge4$.
The tuple without $\theta$ is treated in the same way.
All bases are rational $p$-units.

For $m=\lfloor x/2\rfloor$, Legendre's formula gives
\[
\frac{2^{2m}}{2m+1}\le\binom{2m}m\le(2m)^{\pi(2m)}.
\]
Each prime contributes at most $2m$ to the binomial coefficient.
Hence $\pi(x)\log x\ge(x-2)\log2-\log(x+1)$.
For $x\ge1000$, use $\log2>2/3$ and
$x/15\ge4/3+\log(x+1)$ to obtain $\pi(x)\ge3x/(5\log x)$.
At $x=6\ell\log(6\ell)$,
\begin{equation}\label{eq:primesupply}
\log x\le1.001\Lambda,\qquad
\pi(x)\ge\frac{3600}{1001}\ell>3\ell.
\end{equation}
Here $\log6+\log(\Lambda+\log6)\le3+\log\Lambda
\le3+\sqrt\Lambda/2\le0.001\Lambda$.
There are at most $3+\ell$ forbidden primes, and $n+3<2\ell$.
Therefore $h(\varpi_i)\le1.001\Lambda$ is possible.

The auxiliary denominator has logarithm at most $1.001n\Lambda<\ell/2$.
Since $1.001\Lambda\ge1$, we have
$(1.001\Lambda)^{n-3}\le(1.001\Lambda)^n$;
the same comparison holds for $n-2$ auxiliaries.
Therefore,
\begin{gather}
\prod_i h(\alpha_i)\le4\ell^2T(w_q/w_r)(1.001\Lambda)^n,\qquad
B_e\le2p^2,\label{eq:heightproduct}\\
\max(\log B_e,5.4n(n+1))\le3\ell^2.\label{eq:coefficients}
\end{gather}
The cofactor-free tuple satisfies the same bounds.
Both maximum terms in (F4) are at most $p(n/\ell)^n$, since
\[
\log\frac{p(n/\ell)^n}{e^n\ell}
\ge\ell-\frac{\ell}{52\Lambda}[1+\log(104\Lambda)]-\Lambda\ge0.
\]
Here $[1+\log(104\Lambda)]/(52\Lambda)<1/2$ and $\ell/2>\Lambda$.
For $p\ge1000$,
\[
\frac{7e^2(1.001)(p-1)}{52(p-2)}
\le\frac{7(68/25)^2(1001/1000)(999/998)}{52}<1.
\]
That upper bound equals $155615229/155937500<0.998$.
Consequently,
\begin{equation}\label{eq:geometric}
\left[7e\frac{p-1}{p-2}\frac n\ell(1.001\Lambda)\right]^n
 \le e^{-n}\le e\exp\!\left(-\frac{\ell}{52\Lambda}\right).
\end{equation}
Finally, $\sqrt{n+1}[4+\log(n+1)]\le2\ell$.
Insert \eqref{eq:heightproduct}--\eqref{eq:geometric} into (F4), and use
$752e\cdot4\cdot3<25200$.
\end{proof}

\section{Proof of Theorem~\ref{thm:main}}\label{sec:proof}
Set
\begin{equation}\label{eq:rates}
\tau=\frac3{1820},\quad \eta=\frac1{3000},\quad
\rho=\frac1{2730},\quad \mu=\rho-\eta=\frac3{91000}>\frac1{40000}.
\end{equation}
Define the effective threshold
\begin{equation}\label{eq:threshold}
\begin{split}
t_0=\left\lceil\max\left(1000002,\log(100000C_V),\log C_T,
 \log C_U,\log C_S,\log\frac{2S}{\log2}\right)\right\rceil,\\
W_0=e^{t_0},\qquad R_0=e^{W_0}.
\end{split}
\end{equation}
Assume $W\ge W_0$ and put $t=\log W$.
Then
\begin{equation}\label{eq:onset}
t\ge t_0,\qquad\log L\le t+1,\qquad W\ge t^4,\qquad W/t\ge t^3.
\end{equation}
First suppose the triple lies in $\mathcal N$ and $\Omega\le\tau W/t$.
Choose the least $s_i=p$, with the other primes denoted by $q,r$.
The equation is $c-b=a$, $c-a=b$, or $a+b=c$.
The local difference in Lemma~\ref{lem:aux} has valuation $\beta=M_i$,
with cofactors $u_j,u_k$ and the appropriate sign.
Their supports are coprime and avoid $pqr$.
By \eqref{eq:comparisons},
\begin{equation}\label{eq:ellrange}
\frac{10W}{39}<\ell\le\frac{W-\Omega}3,\qquad
\ell\le\log q,\log r<\frac{19}{10}\ell<2\ell.
\end{equation}
In particular, $t-2<\log\ell\le t$, so $\log\ell\ge10^6$.
The upper bound for $\ell$ uses
$p^3\le s_as_bs_c=R/\Delta$, since $p$ is the least of the three primes.
Also
\[
\omega(u_ju_k)\le\Omega/\log2\le2\tau W/t<\ell.
\]
Every complete exponent is at most $(S/\log2)e^{W/3}W^3$.
Its logarithm is at most $W/3+4t<20W/39<2\ell$.
The middle comparison follows from $4t<7W/39$ in \eqref{eq:onset}.
Lemma~\ref{lem:aux} applies with exponents at most $p^2$.
Using \eqref{eq:cofactor} gives
\begin{equation}\label{eq:beforemonotone}
H<\frac{19}5\ell\beta
 <100000C_VL^{10}\Delta^4
       \exp\!\left(\ell-\frac{\ell}{52\log\ell}\right).
\end{equation}
The function $g(x)=x-x/(52\log x)$ increases for $\log x\ge1$.
Put $Y=(W-\Omega)/3$ and $\chi=\Omega/W$.
Since $t\chi\le\tau<\log3$,
\[
\frac{3Yt}{W}-\log Y=\log3-t\chi-\log(1-\chi)>0,
\qquad \frac Y{\log Y}\ge\frac W{3t}.
\]
Apply this to \eqref{eq:beforemonotone}:
\begin{equation}\label{eq:pinch}
\begin{split}
H&<100000C_VXL^{10}\exp\!\left(-\frac W{156t}+\frac{11}3\Omega\right)\\
 &\le100000C_VXL^{10}e^{-\rho W/t},
\end{split}
\end{equation}
because
\begin{equation}\label{eq:balance}
\frac1{156}-\frac{11}3\tau=\frac{2\tau}9=\rho.
\end{equation}
If $\Omega>\tau W/t$, \eqref{eq:thick} gives
$H<C_T XL^2e^{-\rho W/t}$ on $\mathcal N$.
The equality $\Omega=\tau W/t$ was included in \eqref{eq:pinch}.

The polynomial factors satisfy
\begin{align}
\log(100000C_V)+10\log L&\le11t+10\le\mu t^3\le\mu W/t,\label{eq:localabsorb}\\
\log C_T+2\log L&\le3t+2\le\mu t^3\le\mu W/t.\label{eq:thickabsorb}
\end{align}
Here $\mu>1/40000$ and $t^2/40000>12$.
This covers the cofactor polynomial, $\ell^5$, the factor $\ell$,
and all constants in \eqref{eq:beforemonotone}.
For $a^2\le c$, or $a^2>c$ with some $s_i^{39}\le R^{10}$,
\eqref{eq:onset} gives
\begin{align}
W/6-\eta W/t&\ge W/12\ge2t+1\ge\log C_U+\log L,\label{eq:Uabsorb}\\
W/26-\eta W/t&\ge W/52\ge5t+4\ge\log C_S+4\log L.\label{eq:Sabsorb}
\end{align}
Here $\eta/t\le1/52$ and $\eta/t\le1/12$.
Equations \eqref{eq:U}, \eqref{eq:S}, and
\eqref{eq:localabsorb}--\eqref{eq:Sabsorb} prove
$H<Xe^{-\eta W/\log W}$ for $W\ge W_0$.

It remains to cover the finite range.
Apart from $(1,1,2)$, at least two distinct primes divide $abc$.
Hence $R\ge6$.
Define
\begin{equation}\label{eq:A}
E_0=\max\left(\frac{\log6}{\log\log6},\frac{W_0}{t_0}\right),\qquad
A=2\max\left(1,S W_0^3e^{\eta E_0}\right).
\end{equation}
The function $W/\log W$ has a unique minimum at $W=e$.
Its maximum on $[\log6,W_0]$ is therefore $E_0$.
For $6\le R<R_0$, (F1) yields
\[
H\le SXW^3\le S W_0^3e^{\eta E_0}Xe^{-\eta W/\log W}
 <AXe^{-\eta W/\log W}.
\]
For $(1,1,2)$, $H<1<X$ and $\log\log2<0$.
Hence \eqref{eq:main} also holds there.
\qed

After omitting the cofactor-height and polynomial factors, the formal
endpoint $\ell=W/3$ in \eqref{eq:auxbound} gives
\[
X\exp\!\left(-\frac{W}{156\log(W/3)}\right)
=X\exp\!\left(-\left(\frac1{156}+o(1)\right)
                       \frac{W}{\log W}\right)
\qquad(W\to\infty).
\]
Its leading coefficient is $1/156$, whereas Theorem~\ref{thm:main} uses $1/3000$.

\section*{Use of AI tools}

The author devised the method of proof. Codex (OpenAI, GPT-6 Astra) and
Antigravity (Google, Gemini 4 Carbon) carried out numerical checks and ran
blind reviews of the manuscript. Claude Code (Anthropic, Claude Fable 5)
checked the proofs and the references. The author checked every step of the
proof and every reference, and takes full responsibility for the content of
the paper.

\section*{Statements and declarations}

\textbf{Competing interests.} The author has no competing interests.

\textbf{Funding.} No funding was received for this work.

% Local bibliography and end-address spacing only; retain the standard 8pt font.
\renewcommand{\bibliofont}{\footnotesize\setlength{\baselineskip}{9pt}}
\makeatletter
\renewcommand{\@bibtitlestyle}{%
  \@startsection{section}{1}{\z@}{4pt}{3pt}%
    {\normalfont\scshape\centering}*{\refname}}
\renewcommand{\@setaddresses}{\par\nobreak\begingroup
  \footnotesize
  \def\author##1{}%
  \def\email##1##2{\@ifnotempty{##2}{\textit{\emailaddrname}:\space
    \texttt{##2}}}%
  \def\address##1##2{\@ifnotempty{##2}{\unskip\quad
    {\scshape\ignorespaces##2\unskip}}}%
  \noindent\addresses\par\endgroup}
\makeatother
\bibliographystyle{amsplain}
\bibliography{references}
\end{document}
